\section{Introduction}
\label{sec:introduction}

A manipulation plan is a sequence of contacts. The robot takes hold of the
object, moves it, lets go, and does so again as many times as the task
requires. Choosing that sequence is a combinatorial problem whose cost grows
exponentially with its length, and the length itself is unknown before
planning, since it follows from the goal, from the joint limits of the robot
and from the obstacles around the object. The usual way to attack it is to
build a graph whose nodes gather the configurations in which the robot can
grasp or release the object, and whose edges are the motions between
them~\cite{simeon2004manipulation}, and then to search that graph by sampling
it~\cite{hauser2010multimodal,mirabel2016hpp}, by solving a mixed-integer
programme over it~\cite{deits2014footstep}, or by fixing the sequence of modes
beforehand and optimising within
it~\cite{toussaint2015lgp,levit2025regrasp}.

Each of these needs a model of reachability between contacts: which grasps can
follow the current one, and how far the object can be turned or carried while
it is held. Such a model is nonlinear, because it involves the forward
kinematics of the arm and collision with the environment, and the
configurations that satisfy it occupy thin, lower-dimensional sets. Handling
it exactly at the same time as the combinatorics is impractical, so it is
normally approximated, either by precomputing a finite set of
grasps~\cite{wan2017regrasp,lertkultanon2018certified} or by sampling.
Convex decompositions of the collision-free configuration
space~\cite{deits2015computing,dai2024certified,petersen2023growing,werner2024approximating}
are the continuous alternative: they cover free space with convex sets, over
which Graphs of Convex Sets return trajectories that are globally optimal with
respect to the cover~\cite{marcucci2023motion}.

Convex models of this kind have so far been consumed by mixed-integer
programmes, which must be given the number of grasps and whose indicator
constraints lose their meaning once the sets are thin
(Section~\ref{sec:method:miqp}). Searches that leave the number of grasps open
work instead on a discretisation of the grasps, which both hides solutions and
enlarges the search. This is unfortunate, because the properties we want here
are the ones A* already has: it is deterministic, it recovers the shortest
sequence under an admissible heuristic, and the number of contacts comes out
of the search rather than going into it. The missing piece is a reachability
model that A* can expand without discretising it.

For footstep planning, CASSR~\cite{wang2026cassr} supplies such a model by
expanding convex sets of reachable footholds and fitting the footsteps to the
chosen sets afterwards. Manipulation does not inherit this directly. The
contact sets are not handed over by the environment but have to be constructed
from the grasp constraints, and contact displaces the object, so what is
reachable depends on the configuration of the object as much as on that of the
robot.

We therefore search over sets of grasping configurations instead of individual
grasps. Intersecting a convex decomposition of the collision-free space of
robot and object with a convex model of the grasp constraints yields the
\emph{transition sets}, the places where the robot can take hold of the
object, move it and let go. Our A* expands these sets, and carries with each
one the object configurations it makes reachable. When a single object
coordinate is tied to a single robot joint, as happens for caps, valves and
knobs, those configurations form an interval, which the search propagates and
compares without approximation. Our contributions are:
\begin{itemize}
  \item \emph{transition sets}, a continuous model of grasp reachability.
    They are convex, they are obtained from a convex decomposition of free
    space and a convex model of the grasp constraints, and we validate them
    against a collision checker where the decomposition is least reliable
    (Sections~\ref{sec:problem},~\ref{sec:method:free}--\ref{sec:method:graph});
  \item a \emph{reachability A*} that expands them. It takes no bound on the
    number of grasps, builds its graph from linear programs alone, and returns
    a plan with the fewest grasps its sets allow or establishes that none
    exists (Sections~\ref{sec:method:search}--\ref{sec:method:properties});
  \item a procedure that turns the returned sequence of sets into robot and
    object configurations without an optimiser
    (Section~\ref{sec:method:recovery}), and an evaluation against
    mixed-integer and sampling-based planners up to a 10-dimensional
    robot-object space (Sections~\ref{sec:experiments}--\ref{sec:results}).
    \todo{update with E1--E5.}
\end{itemize}
